\appendix

\section{Counterexample and Corrected Lean Proof (I5a)}
\label{app:counterexample}

\paragraph{Witness instantiation.}
The following concrete values make the original (bound-free)
statement of \texttt{actor\_critic\_two\_timescale\_converges} false:
feature map $\psi(s, a, i) = \uparrow s$ (coercion $\mathbb{N} \to
\mathbb{R}$; unbounded in $s$), step sizes $\alpha_w(t) = 1/(t+1)$
and $\alpha_\theta(t) = 1/((t+1)(1+\log(t+2)))$ (both Robbins-Monro),
trajectory $\tau.\text{states}(t) = t$ (fresh state each step),
rewards $\equiv 1$. Because each state is visited exactly once, the
TD error is $\delta_k = 1$ for all $k$, and the actor accumulates
$\theta_t(0) = \sum_{k=1}^{t-1} k/((k+1)(1+\log(k+2))) \to \infty$.

\paragraph{Corrected statement.}
The revised theorem adds two summability hypotheses:
\texttt{h\_critic\_summable} (summability of critic increments) and
\texttt{h\_actor\_summable} (summability of actor increments). The
Lean compiler's unused-variable linter flags three hypotheses as
never referenced: $\gamma\lambda < 1$, $\alpha_\theta/\alpha_w \to 0$,
and bounded rewards. All three are sufficient conditions that imply
summability. None is necessary. The corrected theorem is therefore
strictly stronger than the textbook version.

\paragraph{Machine-check status.}
\texttt{lake build} exits cleanly, zero \texttt{sorry}, axiom audit
$\{\texttt{propext}, \texttt{Classical.choice}, \texttt{Quot.sound}\}$.
The full Lean source is in \texttt{lean/CreditAssignment/ActorCritic.lean}
in the replication package.

\section{Lean Proof Inventory}
\label{app:lean-inventory}



The three Robbins-Monro convergence theorems (I1a, I2a, I3a) are
proved modulo a shared increment-summability hypothesis that encodes
the stochastic-approximation theorem, pending upstream Mathlib
formalization. Additionally, the Vogels-Sprekeler convergence and
stability theorems (J1a, J1c) were found to be \emph{false as stated}
and corrected. The original actor-critic ablation target (I5a) was
also found false. The corrected version is proved.

\section{Simulation Parameters}
\label{app:sim-params}

All simulations use the Brian~2 spiking neural network simulator
\citep{stimberg2019brian2}. The paradigm follows \citet{tang2024dynamic}:
closed-loop optogenetic stimulation, 100\,ms time step, 55 actions
per session. The eligibility-trace width $\lambda$ is swept over
$[0.7, 0.99]$ in steps of 0.01, with 100 seeds per $\lambda$ value.
The annealing RBF kernel uses $\sigma_0 = 0.5$\,s with exponential
decay $\tau_\sigma = 10$ sessions. Full parameter tables are in the
replication package under \texttt{data/sim\_params.json}.

\section{Mutation Catalog}
\label{app:mutations}

Ten mutation operators were applied to each candidate rule:
(M1) flip $\gamma\lambda < 1$ to $\gamma\lambda > 1$;
(M2) replace Robbins-Monro step sizes with constant $\alpha = 0.1$;
(M3) remove the eligibility trace ($\lambda = 0$);
(M4) replace the RBF kernel with a flat kernel;
(M5) swap actor and critic step sizes;
(M6) remove the bounded-reward hypothesis;
(M7) replace summability with $\ell^2$ condition;
(M8) flip the credit-window direction;
(M9) replace $\gamma = 0.99$ with $\gamma = 0.5$;
(M10) remove the two-timescale condition.
Each mutation was submitted to all six verification modalities.
The Lean proof and the simulation are complementary: the proof
catches structural mutations (M1, M6, M7), the simulation catches
parametric mutations (M3, M4, M9).
